\chapter{Conclusão}\label{cap:conclusoes}

\section{Posição e contribuição}

Este trabalho investigou como avaliar sinais monoculares e inerciais antes de usá-los em uma navegação autônoma. A contribuição principal não é um novo controlador ou um novo algoritmo de planejamento. É um procedimento experimental que liga coleta sincronizada, divisão por execução, referências simples, análise dos erros e um portão espacial de segurança.

A base tabular final reuniu 40 execuções completas e 1.234 intervalos. A MLP apresentou baixa variação algorítmica entre sementes no conjunto fixo, mas não superou a referência constante em todos os alvos. IMU e atitude formaram o grupo com maior efeito médio agregado na ablação, embora o efeito variasse entre alvos; mudanças extremas permaneceram difíceis. Na etapa espacial, a profundidade do Gazebo permitiu completar dez de dez execuções sem adotar caminhos inseguros. O mesmo não ocorreu com as versões monoculares v12 a v19. Elas podiam obter métricas por pixel aceitáveis e, ainda assim, produzir falso espaço livre ou poucos caminhos executáveis.

Esses resultados sustentam a posição defendida: antes de aumentar a complexidade da arquitetura ou colocar o estimador no laço de voo, é preciso verificar a representação e a consequência espacial dos erros. O bloqueio do SITL monocular não é uma ausência de resultado. Ele mostra que a métrica por pixel, isolada, seria um critério insuficiente de liberação.

\section{Respostas às perguntas de pesquisa}

O aumento do conjunto trouxe ganhos modestos e irregulares, embora cinco dos seis alvos tenham terminado com erro menor. As cinco sementes produziram pouca variação no MAE, indicando estabilidade algorítmica no conjunto fixo. IMU e atitude tiveram a maior contribuição agregada entre gradientes e ablação, com variação entre alvos. Informar o evento real ao regressor produziu apenas uma melhora pequena, portanto o erro ficou concentrado na estimação do sinal e da magnitude. Por fim, métricas aceitáveis por pixel não garantiram utilidade espacial, pois nenhuma versão monocular aprovou o portão completo.

\begin{table}[htb]
\centering
\caption{Perguntas de pesquisa, experimentos e respostas.}
\label{tab:mono-perguntas}
\resizebox{\textwidth}{!}{%
\begin{tabular}{p{4.3cm}p{4.2cm}p{6.0cm}}
\hline
\textbf{Pergunta} & \textbf{Experimento} & \textbf{Resposta} \\
\hline
Mais dados melhoram o desempenho? & Marcos de 9 a 40 execuções & Ganhos modestos, irregulares e dependentes do alvo. \\
O resultado é estável? & MLP com cinco sementes e referências & Baixa variação entre sementes; nenhum modelo venceu todos os alvos. \\
Quais entradas contribuem? & Gradientes e ablação & IMU e atitude tiveram a maior contribuição agregada, com variação entre alvos. \\
Onde se concentra o erro de evento? & Classificador, regressor e evento real & Principalmente na regressão do sinal e da magnitude. \\
Métrica por pixel garante caminho útil? & Repetição espacial e A* & Não; as versões monoculares falharam no portão espacial. \\
\hline
\end{tabular}}
\legend{Fonte: elaborado pelo autor.}
\end{table}

\section{Limitações}

Os experimentos foram realizados em um único cenário simulado, com um modelo de veículo, uma rota principal, obstáculos estáticos, ausência de vento e sem variação controlada de iluminação ou clima. A profundidade de referência do Gazebo não incorpora um modelo realista de ruído de sensor. A faixa vertical estreita foi adequada à rota em altitude quase constante, mas não generaliza automaticamente para subidas, descidas ou obstáculos suspensos.

Os atributos tabulares resumem as imagens e perdem informação espacial. Os eventos de proximidade são desbalanceados, o que limita a interpretação de F1 alto. MLP, árvores, os cinco marcos e as ablações foram repetidos com cinco sementes. Os IC95, testes pareados e tamanhos de efeito reduzem a assimetria anterior, mas cinco pares ainda oferecem pouca potência, e as comparações exploratórias entre seis alvos não receberam correção de multiplicidade.

Na validação monocular, o conjunto primário da v19 também participou da validação durante o treinamento e não é independente do desenvolvimento. Os quatro conjuntos secundários foram processados depois do congelamento, mas permanecem no cenário \texttt{baylands}; não houve avaliação fora dele. A reprovação offline impediu o voo SITL monocular, portanto o comportamento integrado do estimador em voo não foi medido. Também não foi usado um relógio comum desde o recebimento do RGB, passando por pré-processamento, inferência, integração e A*, até a publicação; os tempos parciais não constituem latência ponta a ponta.

O código, as configurações, os resultados derivados e os hashes estão versionados, mas os dados brutos não foram publicados devido ao volume. A reprodução integral por terceiros depende de acesso à base do laboratório. O script \path{scripts/generate_tabular_tcc_figures.py} regenera as três figuras tabulares de resultados. As quatro figuras metodológicas antigas ainda não possuem um comando autônomo de regeneração.

\section{Agenda de trabalhos futuros}

O passo seguinte é ampliar a diversidade dos pares RGB--profundidade, incluindo outras rotas, alturas, iluminação e cenas, e reservar desde o início um conjunto espacial que não participe de treinamento nem de escolha de hiperparâmetros. Também é necessário calibrar a incerteza da profundidade por região da imagem e incorporá-la à ocupação, em vez de usar somente margens globais.

Depois da aprovação offline, o modelo deve ser testado primeiro em missões curtas e supervisionadas, mantendo o mapa de referência como guardião. A avaliação em hardware real deve começar com baixa velocidade e ambientes controlados. No eixo tabular, arquiteturas temporais e atributos espaciais podem ser comparados à MLP atual, sempre preservando a divisão por execução e as mesmas referências.
